\section{Cơ sở lý thuyết và công trình liên quan}
\subsection{Hỏi đáp và RAG}
NewsQA yêu cầu tìm đáp án trích từ một bài báo cho sẵn~\cite{newsqa}. NewsLens mở rộng sang hỏi đáp trên toàn kho bằng RAG: truy xuất các đoạn liên quan rồi dùng chúng làm ngữ cảnh cho mô hình sinh~\cite{rag}.
\subsection{Truy xuất và chia đoạn}
\begin{itemize}
\item \textbf{Truy xuất thưa (sparse):} biểu diễn văn bản theo các chiều từ vựng. BM25 dùng tần suất từ, độ hiếm và chuẩn hóa độ dài~\cite{bm25}; learned-sparse của BGE-M3 học trọng số từ vựng~\cite{bge}.
\item \textbf{Truy xuất vector (dense):} ánh xạ câu hỏi và đoạn văn vào vector liên tục, rồi xếp hạng theo độ tương tự.
\item \textbf{Truy xuất kết hợp (hybrid):} kết hợp nhiều danh sách truy xuất. NewsLens dùng nghịch đảo thứ hạng có trọng số (Reciprocal Rank Fusion, RRF)~\cite{rrf}.
\end{itemize}
Bộ xếp hạng lại (reranker) dùng cross-encoder chấm từng cặp câu hỏi--đoạn văn để sắp lại ứng viên; không tìm thêm đoạn ngoài danh sách truy xuất.

Chia đoạn (chunking) tạo đơn vị lập chỉ mục. Cách chia đệ quy (recursive) ưu tiên ranh giới đoạn, dòng và từ trong giới hạn token; phần chồng lấp giữ ngữ cảnh ở biên. Kích thước đoạn, số ứng viên và số đoạn đưa cho mô hình sinh là các tham số riêng.
\subsection{Sinh và đánh giá}
Prompt zero-shot không kèm ví dụ; one-shot có một ví dụ minh họa. Chỉ dấu \texttt{[n]} trong câu trả lời trỏ tới đoạn ngữ cảnh thứ $n$, khác với trích dẫn học thuật trong báo cáo.

NewsLens kết hợp so khớp đáp án, đối chiếu nhãn bằng chứng và điểm chấm bằng mô hình ngôn ngữ của RAGAS~\cite{ragas}. Các thước đo ở Chương 5 đánh giá riêng độ đúng, mức hỗ trợ của ngữ cảnh và dẫn nguồn.
